% ============================================================
%  USPSC-Cap3-Metodologia_3.3-Rotulagem_Manual.tex
%  Seção 3.3 — Rotulagem Manual do Corpus
% ============================================================

\section{Rotulagem Manual do Corpus}
\label{sec:rotulagem}

A rotulagem manual constitui etapa fundamental para a construção do conjunto de referência (\textit{gold standard}) utilizado na avaliação das abordagens de classificação de sentimento descritas na \autoref{sec:abordagens}. O processo foi conduzido pelo próprio pesquisador por meio de uma ferramenta de anotação desenvolvida especificamente para este trabalho, implementada em Python com a biblioteca \texttt{streamlit} e integrada ao banco de dados PostgreSQL descrito na \autoref{subsubsec:pipeline_coleta}.

\subsection{Protocolo de Anotação em Duas Etapas}
\label{subsec:protocolo_anotacao}

O protocolo de anotação foi estruturado em duas etapas sequenciais, aplicadas a cada publicação individualmente:

\textbf{Etapa 1 — Relevância financeira.} O anotador classifica a publicação em uma de duas categorias mutuamente exclusivas: \textit{tweet financeiro} (valor~1), quando o conteúdo se relaciona diretamente ao mercado financeiro, à economia, a investimentos ou a eventos com impacto sobre ativos financeiros; ou \textit{não financeiro} (valor~0), quando o conteúdo é de outra natureza — editorial, esportivo, de entretenimento ou sem relação com o contexto econômico-financeiro. Apenas as publicações classificadas como financeiras prosseguem para a segunda etapa.

\textbf{Etapa 2 — Polaridade de sentimento.} Para as publicações identificadas como financeiras, o anotador atribui uma das três classes de polaridade:

\begin{itemize}
	\item \textbf{Positivo}: o conteúdo veicula uma perspectiva favorável ao mercado ou a agentes econômicos — crescimento, resultados acima do esperado, alívio de incertezas, aprovação de
	medidas de estímulo, entre outros;
	
	\item \textbf{Neutro}: o conteúdo tem caráter predominantemente informativo e factual, sem atribuir valoração explícita ao evento relatado;

	\item \textbf{Negativo}: o conteúdo veicula uma perspectiva desfavorável — queda de indicadores, resultados abaixo do esperado, aumento de incerteza, crises ou eventos de risco sistêmico.
\end{itemize}

A separação entre relevância financeira e polaridade de sentimento em etapas distintas é uma decisão metodológica deliberada: ela evita que publicações de conteúdo não financeiro --- mesmo que textualmente negativas ou positivas --- contaminem a análise de sentimento do mercado financeiro.

\subsection{Ferramenta de Anotação}
\label{subsec:ferramenta_anotacao}

A ferramenta de anotação constitui o segundo estágio do \textit{pipeline} (\texttt{app/dashboard/}) e é composta por dois arquivos:

\begin{itemize}
	\item \texttt{app/dashboard/app.py}: \textit{launcher} multi-página Streamlit, invocado via \texttt{streamlit run app/dashboard/app.py} (acessível em \texttt{http://localhost:8501}). Agrega em uma única interface 6 páginas:
	
	\begin{itemize}
		\item \textit{Anotação} (\texttt{pages/annotation.py});
		\item \textit{Eploração de Dados} (\texttt{pages/exploration.py});
		\item \textit{Split do Dadaset} (\texttt{pages/dataset\_split.py});
		\item \textit{Fine-tuning BERTimbau} (\texttt{pages/fine\_tuning.py});
		\item \textit{Classificação} (\texttt{pages/classification.py}) e;
		\item \textit{Avaliação} (\texttt{pages/evaluation.py}.
	\end{itemize}

\end{itemize}

A interface de classificação é composta por três componentes principais:

\begin{enumerate}
	\item \textbf{Exibição do conteúdo}: para cada publicação, a ferramenta exibe o texto integral do campo \texttt{note\_tweet} e a data de publicação, permitindo ao anotador contextualizar
	temporalmente o conteúdo antes de classificá-lo;
	
	\item \textbf{Botões de classificação}: dois botões para a decisão de relevância financeira (Etapa~1) e, condicionalmente, três botões para a atribuição de polaridade (Etapa~2). O estado da classificação corrente é destacado visualmente, permitindo revisão e correção a qualquer momento;
	
	\item \textbf{Registro de justificativa}: um formulário opcional (\textit{dialog} Streamlit) permite que o anotador registre, em texto livre, a justificativa para cada decisão — tanto para a relevância financeira quanto para a polaridade atribuída. Esse registro é armazenado nos campos \texttt{why\_is\_finance\_news} e \texttt{why\_sentiment} da tabela \texttt{tweets\_classification}, enriquecendo o \textit{gold standard} com raciocínio explícito passível de auditoria futura.
\end{enumerate}

A navegação pelo corpus é sequencial, com barra de progresso indicando o percentual de publicações já classificadas. A ferramenta prioriza automaticamente as publicações ainda sem classificação
humana, identificadas pelo campo \texttt{has\_human\_classification}. Ao final de cada publicação, um histórico tabular exibe todas as classificações anteriores registradas em \texttt{tweets\_classification}, com os campos \textit{classificador}, \textit{sentimento}, \textit{relevância financeira} e as respectivas justificativas.

\subsection{Armazenamento das Anotações}
\label{subsec:armazenamento_anotacoes}

Cada decisão de classificação é persistida em duas estruturas complementares do banco de dados:

\begin{itemize}
	\item \textbf{Tabela \texttt{tweets}}: os campos \texttt{is\_finance\_news} e \texttt{sentiment} são atualizados diretamente, tornando a classificação imediatamente disponível para as etapas de análise subsequentes;
	
	\item \textbf{Tabela \texttt{tweets\_classification}}: registra o histórico completo de cada decisão com o identificador do classificador (\texttt{classificator = ``Humano''}), as justificativas textuais e os valores atribuídos, garantindo rastreabilidade e possibilitando análises de concordância em estudos futuros.
\end{itemize}


\subsection{Particionamento do Corpus Rotulado para Treinamento e Avaliação}
\label{subsec:particionamento}

O corpus rotulado manualmente cumpre, neste trabalho, dois propósitos metodológicos distintos: (i) constitui o conjunto de referência (\textit{gold standard}) utilizado na avaliação comparativa das abordagens de classificação descritas na \autoref{sec:abordagens}; e (ii) é a fonte de exemplos supervisionados para o \textit{fine-tuning} de modelos pré-treinados ajustados localmente, como o BERTimbau (\autoref{subsec:bertimbau}). Esse duplo uso impõe uma decisão de protocolo: para que a comparação entre abordagens que viram os rótulos durante o treinamento (e.g., BERTimbau ajustado) e abordagens que não os viram (e.g., FinBERT-PT-BR e métodos léxicos) seja metodologicamente honesta, todas as abordagens devem ser avaliadas sobre o \textbf{mesmo subconjunto de teste}, definido a priori, antes de qualquer treinamento \cite{jurafsky2026,caseli2024cap1}.

O particionamento é gerenciado pela classe \texttt{DatasetSplitRepository} (\texttt{app/shared /db/dataset\_split.py}), que implementa o método \texttt{assign\_split()} e persiste o resultado na tabela \texttt{dataset\_split}. A operação é idempotente: uma vez realizada, é reutilizada por todas as execuções subsequentes do \textit{pipeline}; o parâmetro \texttt{force=True} permite regenerá-la caso o corpus rotulado seja ampliado. A \autoref{tab:dataset_split} descreve o esquema da tabela.

\begin{table}[!htbp]
	\caption{Esquema da tabela \texttt{dataset\_split}}
	\label{tab:dataset_split}
	\centering
	\begin{tabular}{llp{8.2cm}}
		\hline
		\textbf{Coluna} & \textbf{Tipo} & \textbf{Descrição} \\
		\hline
		\texttt{tweet\_id} & INTEGER & Chave estrangeira para \texttt{tweets.id}; único na tabela \\
		\hline
		\texttt{split}     & VARCHAR & \texttt{'train'} ou \texttt{'test'} \\
		\hline
		\texttt{fold}      & SMALLINT & Fold do K-Fold ($1$--$4$) quando \texttt{split='train'}; \texttt{NULL} quando \texttt{split='test'} \\
		\hline
	\end{tabular}
	\legend{Fonte: Elaborado pelo autor.}
\end{table}

O procedimento compreende duas etapas sequenciais, ambas com estratificação pela classe de polaridade e semente aleatória fixa (\texttt{random\_state=42}), que preserva a distribuição original das três classes em cada conjunto — condição particularmente relevante dado o desbalanceamento esperado em favor da classe neutra, característico de publicações jornalísticas de natureza factual \cite{loughran2011}:

\begin{enumerate}
	\item \textbf{Hold-out de teste} ($\approx$~15\%): por meio de \texttt{train\_test\_split} com estratificação pelo sentimento, extrai-se um subconjunto fixo reservado exclusivamente para a avaliação comparativa final. Os registros recebem \texttt{split=`test'} e \texttt{fold=NULL}. Nenhum classificador — incluindo o BERTimbau — tem acesso a esses registros durante o treinamento, garantindo que as estimativas de desempenho reportadas no \autoref{Cap:Ava_Exp} sejam livres de viés de seleção;
	
	\item \textbf{Conjunto de treinamento} ($\approx$~85\%): os registros restantes recebem \texttt{split=`train'} e são subdivididos em $K=4$ folds por validação cruzada estratificada (\texttt{StratifiedK Fold}, \texttt{shuffle=True}, \texttt{random\_state=42}), cada um marcado com \texttt{fold}~$\in \{1, 2, 3, 4\}$. Esse conjunto é utilizado exclusivamente pelo protocolo de \textit{fine-tuning} do BERTimbau, detalhado na \autoref{subsubsec:arquitetura_bertimbau}; os demais classificadores (OpLexicon, SentiLex-PT e FinBERT-PT-BR) não requerem dados de treinamento e, portanto, não dependem dessa subdivisão.
\end{enumerate}

A \autoref{fig:dataset_split} ilustra a estrutura resultante do particionamento.

\begin{figure}[!htb]
	\centering
	\caption{\label{fig:dataset_split}Estrutura do particionamento estratificado do corpus rotulado}
	\resizebox{0.82\textwidth}{!}{%
		\begin{tikzpicture}
			[
				box/.style   = {rectangle, rounded corners=3pt, draw=black,fill=gray!12, align=center, minimum height=1.0cm,font=\small},
				fold/.style  = {rectangle, rounded corners=3pt, draw=black!60, fill=gray!5,  align=center, minimum height=0.9cm, font=\small},
				lbl/.style   = {font=\scriptsize, align=center},
				node distance = 0.35cm and 0.4cm
			]
			
			%% Corpus total
			\node[box, text width=11.5cm, minimum height=1.1cm]
			(corpus){Corpus rotulado (tweets financeiros com anotação humana)};
			
			%% Hold-out
			\node[box, text width=3.0cm, below left=0.9cm and 1.8cm of corpus,
			fill=gray!25] (test)
			{Hold-out\\$\approx$~15\%\\\texttt{split=`test'}};
			
			%% Train pool
			\node[box, text width=7.8cm, below right=0.9cm and -1.0cm of corpus]
			(train) {Conjunto de treinamento~$\approx$~85\%\quad\texttt{split=`train'}};
			
			%% Folds
			\node[fold, text width=1.55cm, below=0.6cm of train, xshift=-2.7cm]
			(f1) {Fold 1\\\texttt{fold=1}};
			\node[fold, text width=1.55cm, right=0.35cm of f1] (f2)
			{Fold 2\\\texttt{fold=2}};
			\node[fold, text width=1.55cm, right=0.35cm of f2] (f3)
			{Fold 3\\\texttt{fold=3}};
			\node[fold, text width=1.55cm, right=0.35cm of f3] (f4)
			{Fold 4\\\texttt{fold=4}};
			
			%% Arrows
			\draw[->, thick] (corpus.south) -- ++(0,-0.4) -| (test.north);
			\draw[->, thick] (corpus.south) -- ++(0,-0.4) -| (train.north);
			\draw[->, thick] (train.south)  -- ++(0,-0.3) -| (f1.north);
			\draw[->, thick] (train.south)  -- ++(0,-0.3) -| (f2.north);
			\draw[->, thick] (train.south)  -- ++(0,-0.3) -| (f3.north);
			\draw[->, thick] (train.south)  -- ++(0,-0.3) -| (f4.north);
			
			%% Labels
			\node[lbl, above right=0.05cm and 0.1cm of test.north east]
			{\texttt{random\_state=42}\\estratificado por sentimento};
			\node[lbl, below=0.15cm of f2.south, xshift=0.95cm]
			{\texttt{StratifiedKFold}$(K=4,\;\texttt{shuffle=True},\;\texttt{random\_state=42})$};
			
		\end{tikzpicture}%
	}
	\legend{Fonte: Elaborado pelo autor.}
\end{figure}

A fixação da semente aleatória e a persistência dos identificadores na tabela \texttt{dataset\_split} asseguram a reprodutibilidade integral do protocolo: qualquer execução subsequente do \textit{pipeline} sobre o mesmo corpus rotulado recupera os mesmos conjuntos, condição necessária para que comparações entre abordagens introduzidas em momentos distintos da pesquisa permaneçam metodologicamente consistentes.


\subsection{Limitações do Processo de Rotulagem}
\label{subsec:lim_rotulagem}

O processo de rotulagem foi conduzido por um único anotador — o próprio pesquisador —, o que inviabiliza o cálculo de métricas de concordância interanotadores (e.g., Kappa de Cohen). A subjetividade inerente à tarefa, em especial para publicações de polaridade ambígua ou contextualmente dependente, representa uma fonte de incerteza não quantificada nos resultados, reconhecida como restrição metodológica deste trabalho.
